\documentclass[12pt,a4paper,draft]{report}
\usepackage[left=1cm, right=1cm, top=2cm, bottom=2cm]{geometry}
\usepackage{lscape} % Paket für Querformat
\usepackage{booktabs} % Für bessere Tabellenlinien
\usepackage[ngerman]{babel}
\usepackage[T1]{fontenc}
\usepackage{colortbl} % Für Hintergrundfarben in Tabellen
\usepackage{array} %für bessere Tabellenformatierung
\usepackage{xcolor}
\usepackage{graphicx}
\usepackage{tabularx}
\usepackage{babel}
\usepackage{chemformula} % Hauptpaket für chemische Formeln

\usepackage[
backend=biber,
style=authoryear,
sorting=ynt
]{biblatex} %Paket für Bibliography, auch damit wir Zotero damit verwenden können

\addbibresource{bibliographic resource} %da muss die .bib datei dnn reingeschrieben werden



\title{Übersicht der angesetzten Kulturen}
\author{Linus Benetz}

\newcommand{\RN}[1]{\uppercase\expandafter{\romannumeral#1}}

\begin{document}
	\maketitle
\section{Methanobrevibacter thaueri CW}

\begin{table}[h]
	\centering
	\begin{tabularx}{\linewidth}{c  c  c  >{\centering\arraybackslash}X}
		\toprule
		\rowcolor[gray]{0.9} Nummmer & Medium & Animpf-Datum &  Anmerkung\\ 
		\midrule
		 \RN{1} & DSM 1523 & 08.04.26 &  genutzt zum Überimpfen am 21.04.26 / 10 ml \textbf{abzentrifugiert} / im Kühlschrank gelagert \\
		 
		 \RN{1}.1 & DSM 1523 + CoM & 21.04.26 & Medium enthält evtl. keine Spurenelement-Lsg. SL10 / gewachsen am 24.04 -> genutzt zum überimpfen / gelagert im Kühlschrank\\
		 
		 \RN{1}.1.A & DSM 1523 + CoM & 24.04.26 & gewachsen am 27.04 -> in Kühlschrank gepackt \\
		 
		 \RN{1}.2 & DSM 1523 & 21.04.26 & Medium enthält evtl. keine Spurenelement-Lsg. SL10 / gewachsen am 27.04 -> genutzt zum animpfen -> dann Kühlschrank \\
		 
		 \RN{1}.2.A & DSM 1523 + CoM &27.04.26& gewachssen am 04.05.26 -> überimpft am 4.5. und in Kühlschrank\\
		 
		 \RN{1}.2.A.1 & DSM 1523 + CoM &04.05.26& gewachsen am 06.05.26\\
		 
		 \RN{1}.2.B & DSM 119 + CoM & 24.04.26 & gewachsen 30.04.26 -> unverändert dick am 6.05.26 \\
		 
		 \RN{1}.3 & DSM 119 & 21.04.26 & Medium enthält keine Spurenelement-Lsg SL10 / gewachsen am 23.04.24 -> genutzt für \textbf{\ch{EtOH}-Versuch / gelagert im Kühlschrank}\\
		 
		 \RN{1}.4& DSM 119 + CoM & 21.04.26 & Medium enthält keine Spurenelement-Lsg SL10 / gewachsen am 24.04.24 -> in Kühlschrank gepackt / genutz für \textbf{mikroskopieren}\\ \hline
		 
		 \RN{2} & DSM 1523 & 08.04.26 & ganz leichte Trübung am 17.04. -> deutliche Trübung am 20.04. -> am 21.04 bereits wieder lysiert -> nicht genutzt zum überimpfen / am 24.04. zum \textbf{autoklavieren} gegeben\\ \hline
		 
		 \RN{3} & DSM 1523 & 08.04.26 & bisher nicht gewachsen (23.04.26) -> am24.04. zum \textbf{autoklavieren} gegeben\\ \hline
		 
		 \RN{4} & DSM 1523 + CoM & 17.04.26 & bisher nicht gewachsen (23.04.26) -> zum \textbf{autoklavieren} (4.5.26)\\ \hline
		 
		 \RN{5} & DSM 1523 & 17.04.26 & Headspace beim beimpfen kurz atmosphäre ausgesetzt, wurde nochmal ausgetauscht / bisher nicht gewachsen (23.04.26) -> zum \textbf{autoklavieren} gestellt (29.04)\\ \hline
		 
		 \RN{6} & DSM 119 RF + CoM & 17.04.26 & leichte Schlieren am  23.04.26 / unverändert am 27.04 / unverrändert am 6.5.26\\ \hline
		 
		 \RN{7} & DSM 119 + CoM & 17.04.26 &  gewachsen am 23.04 (Schlieren)-> weiterhin bei 37 °C kultiviert -> schauen wie lange bis lysiert -> am 29.04 genutzt zum \textbf{ausplatieren} -> dann Kühlschrank \\
		 
		 \RN{7}.1 & DSM 119 + CoM & 29.04.26 & gewachsen am 04.05.26\\ 
		 
		 \bottomrule
			\end{tabularx}
			\caption{\textit{Methanobrevibacter} thaueri CW in verschiedenen 20 ml Kulturen - DSM 1523, DSM 119 und DSM 119 RF (Rumenfluid anstatt Sludge) - bei 37 °C und Schütteln kultiviert. Als Headspace wurde immer eine Mischung aus \ch{H2} und \ch{CO2} verwendet. Das Medium wurde immer mit 2 ml Ausgangskultur angeimpft. Als Orginalausgangskultur wurde eine von Max verwendet, ansonsten siehe aufeinander aufbauende Nummerierung.}
			\label{Tab:Mthau}
\end{table}




\section{Methanobrevibacter wolinii SH}

\begin{table}[h]
	\centering
	\begin{tabularx}{\linewidth}{c  c  c  >{\centering\arraybackslash}X}
		\toprule
		\rowcolor[gray]{0.9} Nummmer & Medium & Animpf-Datum &  Anmerkung\\ 
		\midrule
		\RN{1}& DSM 1523 & 17.04.26& Wächst sehr gut, bereits am 20.04.26 Zellen zu sehen / genutzt zum Überimpfen am 21.04.26 / 10 ml \textbf{abzentrifugiert} / gelagert im Kühlschrank\\
		\RN{1}.1 & DSM 1523 & 21.04.26 & In Medium eventuell keine Spurenelement-Lsg. / gewachsen am 23.04. / genutzt zum überimpfen und für \ch{EtOH}-Versuch am 24.04.26 / im Kühlschrank gelagert / wurde \textbf{mikroskopiert} \\
		\RN{1}.1.A & DSM 1523 + CoM & 24.04.26 & dick gewachsen am 27.04 -> 29.04. genutzt zum \textbf{ausplattieren} und überimpfen\\
		\RN{1}.1.A.1 & DSM 1523 + CoM & 29.04.26 & dick gewachsen am 04.05.26\\
		\RN{1}.2 & DSM 1523 & 21.04.26 & In Medium eventuell keine Spurenelemente / Wachstumszeichen am 23.04 / am 24.04 in Kühlschrank gepackt\\
		\RN{1}.3 & DSM 119 & 21.04.26 & Im Medium keine Spurenelement-Lsg. / Wachstum am 23.04. / am 24.04 in Kühlschrank\\ 
		\bottomrule
	
	\end{tabularx}
	\caption{\textit{Methanobrevibacter} wolinii SH in verschiedenen 20 ml Kulturen - DSM 1523, DSM 119 und DSM 119 RF (Rumenfluid anstatt Sludge) - bei 37 °C und Schütteln kultiviert. Als Headspace wurde immer eine Mischung aus \ch{H2} und \ch{CO2} verwendet. Das Medium wurde immer mit 2 ml Ausgangskultur angeimpft. Als Orginalausgangskultur wurde eine von Max verwendet, ansonsten siehe aufeinanderaufbauende Nummerierung.}
	\label{Tab:Mwoli}
\end{table}

\section{\textit{Methanobrevibacter} boviskoreani AbM4}
\begin{table}[h]
	\centering
	\begin{tabularx}{\linewidth}{c c c >{\centering\arraybackslash}X}
		\toprule
		\rowcolor[gray]{0.9} Nummer & Medium & Animpf-Datum & Anmerkungen\\
		\midrule
			\RN{1}& DSM 1523 & 17.04.26& Wächst gut, bereits am 20.04.26 Zellen zu sehen, bilden jedoch größere Präzipate / genutzt zum Überimpfen am 24.04.26 sowie für die \ch{EtOH}-Versuche und \textbf{mikroskopiert} / gelagert im Kühlschrank\\
			\RN{1}.1 & DSM 1523 + CoM & 24.04.26& gewachsen am 27.04.24, wieder mit großen Präziptaten -> genutzt zum \textbf{ausplattieren}\\
			\RN{1}.1.A & DSM 1523 + CoM & 29.04.26&\\
			\RN{1}.1.B & DSM 119 + CoM & 30.04.26& gewachsen am 04.05.26\\
			\RN{1}.2 & DSM 119 + CoM & 24.04.26& gewachsen am 27.04, deutlich dicker als im 1523 und auch größere Präzipate -> lösen sich nach schütteln jedoch auf -> genutzt zum überimpfen + \textbf{zentrifugieren}\\ 
			\RN{1}.2.A & DSM 119 + CoM& 27.04.26& gewachsen am 29.04.26\\
			\bottomrule
	\end{tabularx}
	\caption{\textit{Methanobrevibacter} boviskoreani AbM4 in verschiedenen 20 ml Kulturen - DSM 1523, DSM 119 und DSM 119 RF (Rumenfluid anstatt Sludge) - bei 37 °C und Schütteln kultiviert. Als Headspace wurde immer eine Mischung aus \ch{H2} und \ch{CO2} verwendet. Das Medium wurde immer mit 2 ml Ausgangskultur angeimpft. Als Orginalausgangskultur wurde eine von Max verwendet, ansonsten siehe aufeinanderaufbauende Nummerierung.}
	\label{Tab:Mbovi}
\end{table}


\section{\ch{EtOH}-Versuchsreihe}
 Versuch basiert auf Paper (\textcolor{red}{Zitation noch einfügen}) für M. boviskoreani.
 
 \begin{landscape}
 	\
 \begin{table}[h]
 	\centering
 	\begin{tabularx}{\linewidth}{l c c c c c c >{\centering\arraybackslash}X}
 		\toprule
 		\rowcolor[gray]{0.9} Organismus & Medium & \ch{EtOH}[10 mM]& \ch{EtOH}[20 mM]& \ch{EtOH}[50 mM]& \ch{EtOH}[100 mM]& \ch{EtOH}[200 mM]& Notizen\\ 
 		\midrule
 		M. woli \RN{1}.1 & DSM 1523 + CoM &&&&&& evtl leichte Präzipate am 27.04 gesehen -> vmtl doch noch keine Zellen 29.04.\\
 		M. bovi \RN{1} & DSM 1523 + CoM &&&&&& bei allen Präzipate gesehen am 27.04 -> nicht dicker geworden, evtl doch keine Zellen 29.04\\
 		M. thau \RN{1}.4 & DSM 119 + CoM &&&&&& evtl leichtere Schlieren am 27.04 gesehen bei 10-50 mM -> immer noch genauso, heißt vmtl doch keine Zellen 29.04.\\
 		\bottomrule
 	\end{tabularx}
 	\caption{Experimentreihe für verschieden \textit{Methanobrevibacter} ob diese \ch{EtOH} als Substrat verwenden können. Es wurden jeweils 2 ml einer vorgewachsenen Kultur über impft (s. \ref{Tab:Mbovi}). Der Headspace der angesetzten Reihe besteht aus einem \ch{CO2}\ch{N2}-Gemisch.}
 \end{table}
\end{landscape}
\end{document}